% LuaLaTeX-ja を用いた卒業論文テンプレート（内容更新版）
\documentclass[a4paper,11pt]{ltjsreport}

%======================= パッケージ =======================
\usepackage{luatexja}
\usepackage{amsmath,amssymb}
\usepackage{bm}
\usepackage{graphicx}
\usepackage{ascmac}
\usepackage[top=20mm,bottom=30mm,left=20mm,right=20mm]{geometry}
\usepackage{multirow}
\usepackage{url}
\usepackage{color}
\renewcommand{\bibname}{参考文献}

%======================= 本文開始 =========================
\begin{document}

%--------------------- 表紙 ---------------------
\begin{titlepage}
  \centering
  ~~\vspace{1cm}
 
  {\Huge 卒　業　研　究　論　文 \par}
  {\Huge （2025年度）\par}

  \vspace{8cm}

  \begin{tabular}{rl}
    {\LARGE 題　　目}　 & {\LARGE ：低コスト移動ロボットにおける} \\
                         & {\LARGE ~~VLMを用いたターゲット追従の} \\
                         & {\LARGE ~~性能評価} \\ \\ \\
    {\LARGE 指導教員}　 & {\LARGE ：保坂 忠明} \\　 \\ \\
    {\LARGE 学生番号}　 & {\LARGE ：1512221209} \\　 \\ \\
    {\LARGE 氏　　名}　 & {\LARGE ：小坂 尚璃}
  \end{tabular}

  \vfil
  \begin{flushright}
    明治大学理工学部 電気電子生命学科
  \end{flushright}
\end{titlepage}

\tableofcontents

%===========================================================
\chapter{序論}
\label{chap:intro}
\section{研究背景}
近年，画像と言語を同時に扱う基盤モデル（Vision-Language Model, VLM）が普及し，ロボットの行動決定に直接組み込む研究が活発化している。
その研究動向の中心にあるのは，(1) Web 画像やキャプションから学習した一般知識をもとに未知物体や抽象的な指示にも対処できるようにすること，および (2) 環境内の物体やランドマークをクラスラベルなどの意味情報とともに表現した地図を構築し，その上でナビゲーションを行うことである。
これらのアプローチでは，多様なセンサ情報と高い計算資源を前提とし，豊富な事前学習知識を活かした柔軟なナビゲーション性能が報告されている。
一方，教育・ホビー用途や簡易な現場計測のために用いられる低コスト移動ロボットでは，専用 GPU や高精度センサを搭載することは現実的でない。このような場面では，Raspberry Pi 等の小型コンピュータと単一カメラのみを搭載した機体から画像を取得し，外部 PC やクラウド上の VLM に推論をオフロードして進行方向を決定する構成が有力な選択肢となる。地図を作成・保持・参照せず，単一フレーム画像から直接「次の行動」を決めるため，セットアップが短時間で済み，環境変更時のリセット負荷も小さいという利点がある。
しかし，このような地図非依存・単眼カメラ依存の構成では，視覚ノイズや照度変化により判断が乱れやすく，VLM 推論に数秒以上を要する場合には遅延中に機体が進み過ぎる問題が生じる。高機能な VLM を用いても，この種の構成でどの程度安定した追従動作が得られるか，照度や初期位置ずれ，推論遅延が走行安定性にどのような影響を与えるかについては，体系的な実験報告が十分とはいえない。
低コスト移動ロボットにとっては，このシンプルな構成が現実的な選択肢であり，その性能限界と典型的な失敗パターンを実走行で把握することは，教育・プロトタイピングにおける VLM 活用の実装指針として重要である。

\section{先行研究と課題}
本節では，VLM を用いたロボット制御に関する先行研究を概観し，本研究との相違点と課題を整理する。

近年の代表的な研究として，RT-2\cite{rt2} および VL-Maps\cite{vlmaps} が挙げられる。RT-2 は，Web 上の画像およびテキストで事前学習したモデルをロボットデータで追加学習し，カメラ画像とテキスト指示を入力として，ロボットアームが「どの物体を，どの順番で，どのように動かすか」といった操作手順を直接出力する手法である。これにより，未知物体や多様な表現を含む指示に対しても，把持や配置といった操作をゼロショットで生成可能であることが示されている。ただし，GPU を搭載した計算機および高精度センサを備えたロボットプラットフォームを前提としている。

VL-Maps は，移動ロボットが走行中に取得したカメラ画像から，「どこに」「何が」存在するかという意味情報を環境地図上に逐次書き込むことで，意味ラベル付き地図を構築する手法である。利用者が「椅子の近く」「ソファとテレビのあいだ」などの文章で目的地を指定すると，この地図を検索して適切な位置を推定し，そこまでの経路を計画する。あらかじめ物体位置を教えなくても言語による柔軟な目的地指定が可能である一方，地図の構築と検索には高い計算能力とメモリ資源を要する。

これらのアプローチは，未知物体や新しい言い回しにも対応できる汎用性を有するが，いずれも GPU 搭載マシンや複数センサを備えた高価なロボットを前提としており，教育・ホビー用途の低価格ロボットにそのまま適用することは難しい。

これに対して本研究は，低価格の移動ロボットを対象とし，搭載計算資源やセンサ構成が大きく制約された環境を想定する。そのため，単眼カメラで撮影した画像を外部 PC 上の VLM に送信し，得られた出力から「前進」「左」「右」といった単純なコマンドを生成してロボットを制御する，地図非依存の簡潔な構成を採用する。

この構成には，部品コストや準備時間を抑えられるという利点がある一方で，いくつかの問題が想定される。第一に，VLM の推論に数秒を要する場合，その間もロボットは走行を続けるため，指令と実際の位置のずれが生じやすい。第二に，照明条件やカメラの揺れによって画像が変化すると，VLM の出力が不安定になり，ターゲット追従が破綻する可能性がある。第三に，初期位置がわずかに異なるだけでも，遅延や認識誤りが累積し，軌道の分岐を引き起こすことが考えられる。

高性能ロボットを前提とする先行研究では，これらの要因は主たる検討対象とはなってこなかった。しかし，低コスト移動ロボットで VLM を実用的に活用するためには，上記のような要因が走行安定性に与える影響を実機実験により定量的に把握することが重要である。本研究はこの点に着目し，単眼カメラと外部 VLM のみを用いた最小構成においてターゲット追従実験を行い，遅延・照度・初期位置ずれが走行結果に与える影響を明らかにすることを目的とする。

\section{研究目的}
本研究の目的は，教育・ホビー向けに現実的な「地図なし・単眼カメラ・外部 VLM」構成において，ターゲット追従がどの程度安定して動作しうるかを実走行で評価し，遅延・照度・初期位置ずれが成功率や到達時間，コマンド反転といった走行特性に与える影響を明らかにすることである。
具体的には，OSOYOO 社製 Raspberry Pi Robot Car で取得した画像を PC 側で推論し，視覚判断のみで「黄色いターゲットボックス（以下，ターゲット）に接近し，ターゲットが画像内で十分大きく写った時点で停止する」タスクを遂行させる。明るい環境（照度約 900 lx）での走行セットを基準とし，低照度環境（照度約 100 lx および約 20 lx）および初期位置ずれ条件での試行を追加する。ログと取得画像から到達可否，所要時間，コマンド反転，推論遅延を取得し，教育・プロトタイピングで再現しやすい実験条件と結果の傾向として整理する。
実施範囲が限定されているため，網羅的な定量化には追加実験が必要であり，本稿では実施した条件における走行特性と，観察された範囲での失敗例について報告する。

\section{論文構成}
本論文の構成を示す。第\ref{chap:method}章でハードウェアと制御手順を述べる。第\ref{chap:experiments}章で実験環境と計測項目を説明し，第\ref{chap:results}章で暫定結果を示す。第\ref{chap:discussion}章で考察を行い，第\ref{chap:conclusion}章で結論と今後の課題を述べる。

%===========================================================
\chapter{手法}
\label{chap:method}

\section{ハードウェア構成}
対象機体は OSOYOO社製 Raspberry Pi Robot Car（2023版）である。計算機として Raspberry Pi 5（8GB RAM）を搭載し，広角 USB カメラ，モータドライバ（L298N），ステアリング用サーボから構成する。カメラ解像度は 640×480，フレームレートは 5 fps とし，走行は室内の平坦な床面（長さ約 2.5 m の直線区間）で行う。

\section{通信とソフトウェア構成}
Raspberry Pi 上で `raspi\_agent.py` を 8080 番ポートで起動し，`GET /frame.jpg` と `/stream.mjpg` によるフレーム配信，`POST /command` による `FORWARD/LEFT/RIGHT/BACK/STOP` コマンド受信を提供する。PC 側は同一ネットワークから HTTP でアクセスし，VLM の推論結果に基づきコマンドを送信する。

\section{使用モデルとプロンプト}
機体から取得した画像を PC 側で推論する際に使用する VLM として GPT-5-mini-20250807（ローカル実行）を採用する。単一フレームを入力し，プロンプトで以下を明示する。
\begin{itemize}
  \item 黄色いターゲットボックスを画像から検出し，位置と相対距離を判断する。
  \item 応答は `LEFT`，`RIGHT`，`FORWARD`，`FORWARD\_SLOW`，`BACK`，`STOP` のいずれか一語に限定する。
  \item 中央付近（画像中央から水平±10\%以内）でターゲットが十分大きい場合は `STOP`，中央から外れた場合は左右の差分に従って旋回する。
  \item ターゲットが小さいか遠い場合は `FORWARD`，近づいたら `FORWARD\_SLOW` で減速する。
  \item ターゲットが視野外または判別不能な場合は `BACK` で後退し，視界を確保する。
\end{itemize}
推論に 5--7 秒要するため，ループ間隔は 10 秒とし，各ループで最新フレームを 1 枚取得して命令を更新する。応答は一語のみで，追加説明を含めないようプロンプト内で強制する。

\section{制御手順}
各ループで (1) 直近フレームの取得，(2) VLM への送信と一語応答の取得，(3) 応答をそのままモータ指令に変換して送信，(4) 応答と送信時刻をログ化，の順に実行する。画像取得は 10 秒に 1 回とし，毎ループで最新 1 枚を用いる。速度はモータドライバの PWM を固定し，`FORWARD\_SLOW` のみ半分のデューティ比とする。安全のため，連続して `STOP` が返った場合はモータを無通電にし，外部操作でのみ再開できるようにした。

\section{キャリブレーション}
低コスト機体ではモータやネジの締まり具合，バッテリ電圧により走行量が変動しやすい。実験前に以下の簡易計測を行い，指令と実際の変位を記録した。\footnote{各値は当日の電圧と床面での一回計測を基にしており，追加計測で更新予定。}

\begin{table}[htbp]
  \centering
  \caption{モータ指令と実際の変位（キャリブレーション用）}
  \label{tab:calibration}
  \begin{tabular}{|c|c|c|}
    \hline
    指令 & 変位・回転量 & 備考 \\
    \hline
    FORWARD（1 s） & 計測値を後日挿入 & 電圧と床面を記録 \\
    BACK（1 s） & 計測値を後日挿入 &  \\
    LEFT（0.5 s） & 計測値を後日挿入 & サーボ初期角 90°基準 \\
    RIGHT（0.5 s） & 計測値を後日挿入 &  \\
    \hline
  \end{tabular}
\end{table}

この変位を基に，ターゲットが中央より左に見える場合は `LEFT` を 1 回，右なら `RIGHT` を 1 回送信し，中央付近かつ近距離の場合のみ前進や停止を指示する。変位量は日ごとの再測定で更新し，ログとともに保存する。

%===========================================================
\chapter{実験}
\label{chap:experiments}

\section{環境と照度}
実験は長さ約 2.5 m の室内直線コースで実施した。床は平坦なフロア材で，壁や背景は固定とし，ターゲットは終点に置いた黄色いボックス（“TARGET” と印字）である。初期視野はターゲットが画面中央から水平±10\%，垂直方向は中央よりやや上，サイズは画面比 5--15\% となるように位置決めした。必須セットでは照度計で約 873 lx を維持し，比較のために 300 lx 前後に落とした暗め条件を用意した。

\section{試行セット}
実施した走行セットを表\ref{tab:runs}にまとめる。優先度の高い必須セットを 6 本（時間があれば 10 本）行い，時間が許す場合のみ暗め照度と初期位置ずれの試行を追加した。

\begin{table}[htbp]
  \centering
  \caption{走行セットの構成}
  \label{tab:runs}
  \begin{tabular}{|c|c|c|p{7cm}|}
    \hline
    区分 & 本数 & 照度 & 内容 \\
    \hline
    必須 & 6（推奨 10） & 約 873 lx & 中央±10\% の初期視野，サイズ 5--15\%。同一設定で連続走行し，成功率・到達時間・反転回数・推論遅延を取得。 \\
    オプション1 & 2（時間があれば 4） & 約 300 lx & 初期視野は必須セットと同一。照度だけを下げ，照度依存性を確認。 \\
    オプション2 & 2 & 約 873 lx & 初期視野を左右にずらす（左寄せ -20〜-25\%，右寄せ +20〜+25\%）。水平ずれに対する耐性を確認。 \\
    \hline
  \end{tabular}
\end{table}

実行順は (1) 必須セット 6 本，(2) 暗め 2 本（時間があれば），(3) 左右ずれ 2 本とした。各試行でフレーム画像，VLM 応答，送信コマンド，タイムスタンプ，別カメラで取得した俯瞰画像を保存し，通し番号と照度，開始・終了時刻をメタデータとして付与した。

\section{計測項目}
本稿で報告する主な指標は以下のとおりである。
\begin{itemize}
  \item 成功率：ターゲットが画像内で十分大きく写り，VLM が STOP を出力して機体が静止した試行の割合。
  \item 到達時間：開始コマンド送信から停止判定までの時間。
  \item コマンド反転回数：左右または前後の即時反転の回数。
  \item 推論遅延：VLM 応答時間の中央値と 95 パーセンタイル。
\end{itemize}
コマンド適合率や軌跡偏差など詳細な指標は，ログの追加解析後に必要に応じて追記する。

%===========================================================
\chapter{結果（暫定）}
\label{chap:results}
直近のログには各試行の到達可否，送信コマンド列，推論遅延，照度，初期視野が記録されている。集計は進行中であり，本稿では試行セットごとの代表値を表\ref{tab:summary}に占位した。値が確定し次第，成功率と到達時間中央値，反転回数中央値，推論遅延中央値／95 パーセンタイルを順に埋める。軌跡オーバーレイとコマンド時系列は図\ref{fig:trajectory1}，図\ref{fig:trajectory2}に差し替える予定である。

% 予備の表・図枠（後で数値を埋める）
\begin{table}[htbp]
  \centering
  \caption{試行セット別の暫定集計（後日更新）}
  \label{tab:summary}
  \begin{tabular}{|c|c|c|c|c|}
    \hline
    セット & 照度[lx] & 到達時間中央値[s] & 成功率 & 反転中央値 / 推論遅延中央値・95p \\
    \hline
    必須（6--10 本） & 約 873 & -- & -- & -- / --・-- \\
    暗め（2--4 本） & 約 300 & -- & -- & -- / --・-- \\
    左右ずれ（2 本） & 約 873 & -- & -- & -- / --・-- \\
    \hline
  \end{tabular}
\end{table}

\begin{figure}[htbp]
  \centering
  \fbox{\rule{0.65\textwidth}{0.32\textwidth}}
  \caption{軌跡オーバーレイ（後日実データに置換）}
  \label{fig:trajectory1}
\end{figure}

\begin{figure}[htbp]
  \centering
  \fbox{\rule{0.65\textwidth}{0.32\textwidth}}
  \caption{コマンド系列と推論遅延の時系列（後日実データに置換）}
  \label{fig:trajectory2}
\end{figure}

%===========================================================
\chapter{考察（予定）}
\label{chap:discussion}
現時点のログでは，推論遅延が 1 ループ 10 秒の大半を占めるケースがあり，遅延中に機体が進み過ぎてターゲットを通過する失敗が見られた。照度を 300 lx 前後まで下げた場合，ターゲットの輪郭が弱まり `LEFT`/`RIGHT` の揺れが増える傾向がある。初期位置を左右に 20--25\% ずらすと，旋回 1 回では中心に戻り切らず，反転回数が増えることが予備観察から示唆された。

このため，(1) 推論遅延と移動距離の関係をキャリブレーション表\ref{tab:calibration}と組み合わせて補正すること，(2) 低照度での画像前処理（ガンマ補正など）を入れて VLM 入力を安定させること，(3) 左右ずれ時に二連続旋回を許容するなど簡易な補助規則を加えることが今後の改善候補となる。ログの集計が揃い次第，これらの仮説を定量的に検証する。

%===========================================================
\chapter{結論}
\label{chap:conclusion}
本研究は，低コスト移動体で地図や複雑なセンサ統合を用いない構成において，VLM によるターゲット追従の実走行評価枠組みを示した。OSOYOO社製 Raspberry Pi Robot Car の画像を PC 側で推論し，単一フレームに対する一語応答をそのまま指令化する手順を設計し，照度 873 lx を基準とした走行セットを実施した。記録したログから成功率，所要時間，反転回数，推論遅延を算出し，照度低下や初期位置ずれの影響を追加試行で確認する計画である。今後，キャリブレーションと前処理を組み込み，遅延と走行距離の関係を補正した上で，RT-2 や VL-Maps の知見を参考にした拡張（行動表現の工夫や視覚特徴の共有）の適用可能性を検討する。

%===========================================================
\chapter*{謝辞}
\addcontentsline{toc}{chapter}{謝辞}
本研究の遂行にあたり，実験環境整備に協力いただいた研究室メンバー各位に感謝する。計算資源と機材提供に謝意を表する。

%===========================================================
\bibliographystyle{junsrt}
\bibliography{mybibfile}

%===========================================================
\end{document}
